% ID: SEL-D0002-OA-E02
% ESTADO: cotejada. Fuente y transcripción contrastadas de forma independiente; salida revisada.
% Fuente: Selectividad/sel_2019_matematicas_aplicadas/Titular Septiembre/Exámenes Mat Apl CCSS II - A y B.pdf
% SHA-256: 31fb46947e141f5c8da8f7b381e7d27b67fe5a9d1597120454ce66fce0313f1a
% Página: 1; opción A; ejercicio 2.
% Encabezado original: EJERCICIO 2. Solo enunciado; no contiene solución.

El coste de producción de un bien en una fábrica viene dado por
$C(x)=2(2x-1)^2+1$, con $0\leq x\leq 2$, donde $x$ es la cantidad producida
en millones de kilogramos.

\begin{enumerate}[label=\alph*)]
  \item \textbf{(1 punto)} Estudie el crecimiento y decrecimiento de la función
  $C(x)$.
  \item \textbf{(0.75 puntos)} Determine la cantidad a producir para que el
  coste de producción sea mínimo. \textquestiondown{}Cuál es dicho coste?
  \item \textbf{(0.75 puntos)} Realice un esbozo de la gráfica de la función
  $C(x)$.
\end{enumerate}
